\documentclass[11pt]{article}

\usepackage[margin=1.1in]{geometry}
\usepackage{amsmath}
\usepackage{amssymb}
\usepackage{amsthm}
\usepackage{booktabs}
\usepackage{microtype}
\usepackage[round,authoryear]{natbib}
\usepackage[colorlinks=true,allcolors=blue]{hyperref}

\newtheoremstyle{itheads}{\topsep}{\topsep}%
  {\itshape}{}{\itshape}{.}{ }{\thmname{#1}\thmnumber{ #2}%
  \thmnote{ (#3)}}
\theoremstyle{itheads}
\newtheorem{proposition}{Proposition}
\newtheorem{corollary}{Corollary}

\makeatletter
\renewcommand\paragraph{\@startsection{paragraph}{4}{\z@}%
  {3.25ex \@plus1ex \@minus.2ex}{-1em}%
  {\normalfont\normalsize\itshape}}
\makeatother

\title{\textbf{One More Pivot}\\[0.4em]
\large Farsighted Search When Trials Are Scarce}
\author{Peter Cotton}
\date{\today}

\begin{document}
\maketitle

\begin{abstract}
\noindent Firms search for strategy the way Callander's searchers probe
an unknown correlated landscape: each trial is a strategy actually
lived, and the last one is the one committed to. Field experiments show
the search is short --- roughly three quarters of startups pivot zero
times in a year, a fifth exactly once, and six to eleven percent more
than once --- so the empirically relevant problem is not an infinite
horizon truncated but a game of at most two moves, and that game has
recently been solved exactly on the canonical mean-reverting landscape.
This paper translates the solution into operating rules: when a
struggling incumbent should pivot far, how wide to bracket before the
final commitment, the explicit performance ceiling above which the
grass is provably not greener, and the threshold --- strictly below the
performance of positions already held --- at which a disappointing
pivot should be abandoned rather than repaired. The abandonment rule is proved here; the other rules import
from the companion exact solution. The rules are graded
and falsifiable, the interior option is worth about one percent of
expected value at its best, and the same field experiments that
motivate the model rhyme with its central prediction: better search
produces exactly one pivot, not many.
\end{abstract}

\section{Callander's problem}

\citet{callander2011b} posed the problem this way. Agents do not know
how choices map to outcomes; the mapping is the realized path of a
Brownian motion, known where it has been tried, unknown elsewhere, with
nearby choices performing similarly. Search is by trial and error:
each agent observes the choices of earlier agents and the outcomes
they realized, then makes its own choice and lives the result. Trials
are not laboratory queries; each is a product shipped or a position
occupied, and Callander's own application is new product development,
where the model reproduces features of the product life cycle. In the
policymaking companion \citep{callander2011}, the same machinery shows
that trial-and-error policymaking can become stuck at outcomes
arbitrarily bad, with no institutional constraint to blame.

The workhorse of the tradition is myopic by construction, not by
assumption: each agent in the sequence chooses once, for its own
outcome, so nobody in the model plans two trials ahead. The survey of \citet{bardhicallander2026} lists
farsighted search --- a searcher who plans her few remaining trials
jointly --- among the tradition's open methodological problems, and
recent progress has bought tractability with restrictions: contiguous
movement along a vein \citep{urgunyariv2025, wong2025}, or exact
few-point placement under an information objective \citep{bardhi2024}.
The management literature on rugged landscapes \citep{baumann2019}
runs on simulation for the same reason. The entrepreneur inverts
Callander's demographic structure: one agent holds the incumbent
position, plans the remaining trials, and lives with the final
commitment. That is the farsighted case, and the evidence of the next
section says it is short.

\section{The evidence: search is short}

Field evidence indicates that real ventures make very few strategic
pivots. In the randomized trial of \citet{camuffo2020}, 74\% of firms
pivoted zero times over roughly a year, 19\% exactly once, and 7\% more
than once. The four-trial replication \citep{camuffo2024}, with 759
firms, found the same shape: mode zero, 11\% more than once, 6\% more
than twice. Practitioner folklore agrees on the mechanism ---
``the true measure of runway is how many pivots a startup has left''
\citep{ries2011} --- though no peer-reviewed cost-per-pivot measurement
exists, and pivot counts should be read with the construct caution of
\citet{kirtley2023}, who show pivots are often accumulated increments
rather than single jumps.

The treatment effect is the striking part. Teaching firms a scientific
approach to search raised the probability of pivoting exactly once by
8.3 percentage points, left twice unchanged, and \emph{lowered} the
probability of pivoting more than twice by 3.7 points
\citep{camuffo2024}. Better search does not produce more search. It
produces one good pivot.

A searcher holding an incumbent position, one further trial, and a
final commitment --- at most two pivots --- is therefore not a
truncation of the empirically relevant problem. It is the empirically
relevant problem.

\section{The model, in business terms}

The landscape is the exponential of a stationary Ornstein--Uhlenbeck
path: log-performance mean-reverts, so there is no free lunch far from
what is known, and the exponential makes rewards right-skewed, as
venture outcomes are. The firm holds an incumbent position of known
performance, may run one more trial anywhere, and then must commit ---
to the incumbent, to the tested alternative, or to an untested position
between or beyond them. The payoff is the performance of the final
position, not the best ever visited: the strategy committed to is the
one lived with.

Mean reversion is what makes the pivot a trade. On a pure random walk
landscape, moving far from the incumbent keeps its expected performance
while adding variance without bound, and a skew-loving searcher would
walk to the horizon; the model would recommend infinite pivots and
contradict the evidence in the first sentence. Under mean reversion,
distance buys variance but surrenders expected performance toward the
population median, and the exact solution below prices that trade.

The exact solution is developed in a companion paper
\citep{cotton2026grass}; here we state what it says operationally.

\section{Operating rules from the exact solution}

\paragraph{Go far on bad news.} With one trial and then commitment, the
optimal placement has three phases. An incumbent below the median
performance should be abandoned with a distant trial: its information
is not worth anchoring on. A strong incumbent, above an explicit
threshold, should be held without experimentation. In between, the
trial goes a distance set by the mean-reversion scale of the industry,
not by the firm's ambition.

\paragraph{Bracket before committing.} With two trials, the searcher
first establishes a bracket of two known positions and then faces the
signature question: commit to a tested position, or to untested ground
between them? Every interior position matches the expected performance
of an explicit tested-side alternative while carrying strictly more
uncertainty, so the interior is preferred exactly when the payoff is
convex --- when upside is what is being bought. For a risk-averse
operator judged on expected utility rather than expected value, the
same identity flips the verdict: the grass between two tested positions
is provably browner.

\paragraph{The ceiling.} Above an explicit performance level, the grass
is not greener anywhere: mean reversion caps the expected performance
of every unvisited position, and no amount of variance-loving
compensates. A firm whose incumbent clears the ceiling has nothing to
gain from further search, and the ceiling is a formula, not a
sentiment.

\paragraph{A bad pivot is abandoned, not repaired.} Entering a bracket
is not a commitment. If the trial between two known positions comes in
at or below the median, every further move inside the spoiled bracket
--- including repairing the disappointing position --- is weakly
dominated by leaving for the best outside option; strictly, if the
draw is strictly below median (Proposition~\ref{prop:exit}, proved in
the appendix). The general rule is a threshold: the value of staying
is nondecreasing in the interior draw while the exit menu does not
depend on it, so the searcher stays exactly when the draw beats a
cutoff (Corollary~\ref{cor:threshold}). In the calibration computed,
the cutoff sits below the performance of the positions already held
--- 0.63 against anchors of 0.70 --- so a pivot that lands somewhat
worse than both of its neighbors is still worth developing, while one
that lands badly is walked away from. The asymmetry is the model's
sharpest managerial statement.

\paragraph{The size of the cleverness.} The interior option --- the
right to commit to untested ground between two tested positions --- is
worth about one percent of expected value at its best. Most of the
value of a second trial is the trial itself, not the sophistication of
where the commitment lands afterward. The model prices its own
subtlety, and prices it small.

\section{The rhyme with the experiments}

The exact solution has a revisit window: an interval of incumbent
performance in which the optimal play is one trial and a considered
commitment, bounded below by the flee region and above by the ceiling.
The randomized trials found that better search shifts firms toward
exactly one pivot and away from both zero and many. The two shapes
rhyme. We state this as a rhyme and not an identity: the treatment
effect in \citet{camuffo2024} is a nonlinearity in an intervention's
effect on pivot counts, not a causal estimate that a third pivot
destroys value, and the model's window is an interval in incumbent
quality, not in pivot count. A design that measured incumbent quality
at baseline could test the correspondence directly.

\section{One last trial in operations}

The same three-move structure appears wherever a process owner holds an
incumbent operating point, has time for one more pilot, and must then
deploy: the final ``shot'' is the deployment decision, not a
laboratory run. The rules above apply verbatim, with one qualification
that decides applicability: the commitment set must genuinely include
untested interpolated settings within an approved operating range. If
deployment is restricted to individually validated settings, the
interior option is off the table and the problem reduces to the
simpler tested-side rules.

For the experimental-design community the model is a benchmark rather
than a rival. The trial-placement step maximizes the expected value of
the best deployable setting after one more observation, which is the
knowledge-gradient criterion of \citet{frazier2008} with an
exponential objective; no new selection principle is claimed. What the
model contributes is an unusually explicit instance --- exact interior
maximization, an exact phase diagram --- against which
knowledge-gradient implementations and their approximations can be
checked, and possibly undercut in computation, since on a Markov
landscape a new observation changes only its containing interval.

\section{A radio interpretation, tested}

The mathematics does not know it is about pivots. A technician
surveying WiFi signal strength holds the same problem: an incumbent
mounting spot with measured power, time for one more measurement, then
the equipment is mounted --- possibly at an unmeasured spot between or
beyond the two readings. Log-power shadowing follows the classical
model of \citet{gudmundson1991}: exponentially correlated in space,
which is exactly the exponentiated Ornstein--Uhlenbeck landscape. On a
public indoor dataset \citep{feng2024} --- a 92-by-15 meter floor, 642
reference points, thirteen access points --- the premise checks
directly: near-Gaussian shadowing residuals and exponential spatial
correlation with a length near seven meters, at $R^2 \approx 0.98$, in
two independent reproductions (\texttt{casestudy/wifi} in the
repository).

The replay is the honest part. On held-out corridors, one extra
measurement lifts the mounted signal by 23\% over the incumbent and
the full three-shot policy by 33\% --- but equispaced far probing and
a knowledge-gradient heuristic do better still, near 70\%. The theory
predicts its own defeat here: the correlation length is about eleven
grid units against corridors of fifteen to ninety, so far probes are
nearly fresh draws, and the small-correlation expansion of the exact
formulas makes the interior option's edge first order in the
correlation while the fresh-draw variance bonus is order one. Where
trials can roam many correlation lengths, measure far and take the
best is leading-order optimal, and the bracket refinements of
Section~4 are second-order corrections.

The business translation is the regime lesson. An industry's
correlation length decides which searcher you are. The radio surveyor
roams tens of correlation lengths and should probe far and simply take
the best. A venture is tethered: capabilities, customers, and identity
keep the next position within roughly one correlation length of the
incumbent, and that is precisely the regime --- near brackets, high
incumbents --- where the phase diagram, the ceiling, and the
abandonment threshold carry the decision. The replay's round-one gaps
are recorded with it: gains are pooled rather than conditional on the
incumbent's standing, per-corridor mean effects are unmodeled, and the
flee move is truncated at corridor ends.

\section{Closing}

Neither the myopic workhorse nor the infinite-horizon limit describes
the search the evidence shows firms conducting. The scarce-trial regime
is its own theory, it now has an exactly solved case, and its rules are
concrete: go far on bad news, bracket before committing, respect the
ceiling, abandon bad pivots at a threshold below your anchors, and
expect the cleverness premium to be small. We leave to future work the $k$-shot game --- in the fast-reversion
regime the data occupies, a singular-perturbation expansion around
independent draws, with the exact three-shot solution as its seed ---
along with information-motivated trials and the field test the rhyme
of Section~5 invites.

\begin{thebibliography}{99}

\bibitem[Bardhi(2024)]{bardhi2024}
Bardhi, A. (2024).
Attributes: selective learning and influence.
\emph{Econometrica}, 92(2):311--353.

\bibitem[Bardhi and Callander(2026)]{bardhicallander2026}
Bardhi, A., and Callander, S. (2026).
Learning in a correlated world.
\emph{Annual Review of Economics}, forthcoming.

\bibitem[Baumann et al.(2019)]{baumann2019}
Baumann, O., Schmidt, J., and Stieglitz, N. (2019).
Effective search in rugged performance landscapes: a review and
outlook.
\emph{Journal of Management}, 45(1):285--318.

\bibitem[Callander(2011a)]{callander2011}
Callander, S. (2011a).
Searching for good policies.
\emph{American Political Science Review}, 105(4):643--662.

\bibitem[Callander(2011b)]{callander2011b}
Callander, S. (2011b).
Searching and learning by trial and error.
\emph{American Economic Review}, 101(6):2277--2308.

\bibitem[Camuffo et al.(2020)]{camuffo2020}
Camuffo, A., Cordova, A., Gambardella, A., and Spina, C. (2020).
A scientific approach to entrepreneurial decision making: evidence
from a randomized control trial.
\emph{Management Science}, 66(2):564--586.

\bibitem[Camuffo et al.(2024)]{camuffo2024}
Camuffo, A., Gambardella, A., Messinese, D., Novelli, E., Paolucci,
E., and Spina, C. (2024).
A scientific approach to entrepreneurial decision making: large scale
replication and extension.
\emph{Strategic Management Journal}, 45(6):1209--1237.

\bibitem[Cotton(2026)]{cotton2026grass}
Cotton, P. (2026).
When the grass is greener: three-shot search on an exponentiated
Ornstein--Uhlenbeck landscape.
Working paper, \texttt{github.com/microprediction/browniansearch}.

\bibitem[Feng et al.(2024)]{feng2024}
Feng, X., Nguyen, K. A., and Luo, Z. (2024).
WiFi RTT RSS dataset for indoor positioning.
Zenodo record 11558192, CC BY 4.0.

\bibitem[Frazier et al.(2008)]{frazier2008}
Frazier, P. I., Powell, W. B., and Dayanik, S. (2008).
A knowledge-gradient policy for sequential information collection.
\emph{SIAM Journal on Control and Optimization}, 47(5):2410--2439.

\bibitem[Gudmundson(1991)]{gudmundson1991}
Gudmundson, M. (1991).
Correlation model for shadow fading in mobile radio systems.
\emph{Electronics Letters}, 27(23):2145--2146.

\bibitem[Kirtley and O'Mahony(2023)]{kirtley2023}
Kirtley, J., and O'Mahony, S. (2023).
What is a pivot? Explaining when and how entrepreneurial firms decide
to make strategic change and pivot.
\emph{Strategic Management Journal}, 44(1):197--230.

\bibitem[Ries(2011)]{ries2011}
Ries, E. (2011).
\emph{The Lean Startup}. Crown Business.

\bibitem[Urgun and Yariv(2025)]{urgunyariv2025}
Urgun, C., and Yariv, L. (2025).
Contiguous search: exploration and ambition on uncharted terrain.
\emph{Journal of Political Economy}, 133.

\bibitem[Wong(2025)]{wong2025}
Wong, Y. F. (2025).
Forward-looking experimentation of correlated alternatives.
\emph{Theoretical Economics}, 20.

\end{thebibliography}

\appendix

\section{The abandonment threshold}

Work on the log scale, where a location with conditional mean $\mu$
and conditional variance $\nu$ has log-value $\mu + \nu/2$, and let
$\zeta(b)$ denote the exterior two-shot value anchored at a known
log-value $b$: explicitly $\zeta(b) = \tfrac12$ for $b < 0$,
$\tfrac12(b^2+1)$ for $0 \le b \le 1$, and $b$ for $b > 1$
\citep{cotton2026grass}. A bracket has known end log-values $b$ and
$d$; an interior trial has revealed log-value $v$.

\begin{proposition}[A bad draw sends you back out]\label{prop:exit}
Let $v \le 0$, the process median or below. Then every placement
interior to the two spoiled sub-brackets, and the revisit of $v$, is
weakly dominated by exiting to $\zeta(\max(b,d))$; strictly if
$v < 0$.
\end{proposition}

\begin{proof}
Consider the sub-bracket with near end $b$ and far end $v$, and an
interior point with correlations $\lambda_2$ and $\lambda$ to its
ends, $\rho_1 = \lambda_2 \lambda$. The conditional mean is
$\mu = A b + B v$ with
$A = \lambda_2 (1-\lambda^2)/(1-\rho_1^2)$ and
$B = \lambda (1-\lambda_2^2)/(1-\rho_1^2)$, both nonnegative, and
$A \le \lambda_2$ because $\lambda \ge \rho_1$ implies
$1-\lambda^2 \le 1-\rho_1^2$. Since $v \le 0$,
$\mu \le \lambda_2 b$. The conditional variance obeys
$\nu = (1-\lambda_2^2)(1-\lambda^2)/(1-\rho_1^2)
\le 1-\lambda_2^2$ by the same inequality. Hence
\[
\mu + \tfrac{\nu}{2} \le \lambda_2 b + \tfrac{1-\lambda_2^2}{2}
\le \sup_{\ell} \Big[ \ell b + \tfrac{1-\ell^2}{2} \Big]
= \zeta(b),
\]
the last identification holding for every real $b$: for
$b \in [0,1]$ the supremum is $(b^2+1)/2$ at $\ell = b$; for $b > 1$
it is $b$ at $\ell = 1$; for $b < 0$ it is $\tfrac12$ at
$\ell = 0$. For $v < 0$ any strictly interior point has $B > 0$, so
$\mu < \lambda_2 b$. The revisit of $v$ pays $v \le 0 < \tfrac12
\le \zeta$. The same bound applies to the other sub-bracket with $d$
in place of $b$.
\end{proof}

\begin{proposition}[Monotone value]\label{prop:monotone}
For the $k$-shot game with any frontier of observed location--value
pairs, the optimal value is nondecreasing in each observed value.
\end{proposition}

\begin{proof}
By induction on $k$. Every candidate placement's conditional mean ---
bridge or exterior ray --- is a nonnegative linear combination of
observed values, and every conditional variance depends on locations
only. At $k=0$ the payoff is a maximum of terms nondecreasing in each
observed value. For the step, raise one observed value and couple the
reveal at any placement upward by the shift in its conditional mean;
the inductive hypothesis and the maximum over placements preserve
monotonicity.
\end{proof}

\begin{corollary}[Threshold rule]\label{cor:threshold}
After an interior reveal $v$ with one shot left, the exit menu is
anchored at outer values and does not depend on $v$, while the stay
menu is nondecreasing in $v$ by Proposition~\ref{prop:monotone}. A
nondecreasing function crosses a constant once: the searcher stays
exactly when $v$ exceeds a threshold $v'(b,d,\rho)$. At
$(b,d,\rho) = (0.7, 0.7, 0.5)$ the threshold computes to $0.6301$,
below both anchors; the computation is
\texttt{verify/check\_commitment.py} in the repository.
\end{corollary}

\end{document}
